\documentclass[12pt, a4paper]{article}
% Shared preamble — included by all Durin's Code LaTeX documents
\usepackage[a4paper, margin=2.5cm]{geometry}
\usepackage{amsmath, amssymb}
\usepackage{graphicx}
\usepackage{xcolor}
\usepackage{listings}
\usepackage{booktabs}
\usepackage{longtable}
\usepackage{array}
\usepackage{enumitem}
\usepackage{titlesec}
\usepackage{fancyhdr}
\usepackage{hyperref}
\usepackage{fontenc}
\usepackage{inputenc}
\usepackage{microtype}
\usepackage{caption}
\usepackage{subcaption}
\usepackage{tikz}
\usepackage{forest}
\usepackage{mdframed}
\usetikzlibrary{arrows.meta, positioning, shapes.geometric, fit, backgrounds}

% ── Colour palette ──────────────────────────────────────────────
\definecolor{durins-blue}{RGB}{31, 78, 121}
\definecolor{durins-gold}{RGB}{180, 130, 20}
\definecolor{durins-gray}{RGB}{90, 90, 90}
\definecolor{code-bg}{RGB}{248, 248, 248}
\definecolor{code-kw}{RGB}{0, 0, 180}
\definecolor{code-str}{RGB}{163, 21, 21}
\definecolor{code-cmt}{RGB}{0, 128, 0}
\definecolor{code-num}{RGB}{9, 134, 88}

% ── Listings style ──────────────────────────────────────────────
\lstdefinelanguage{DurinsCode}{
  morekeywords={room, action, item, npc, exit, if, else, print, remove,
                player, true, false, description, current_room},
  morestring=[b]",
  morecomment=[l]{//},
  sensitive=true
}
\lstdefinestyle{dc}{
  language=DurinsCode,
  backgroundcolor=\color{code-bg},
  basicstyle=\ttfamily\small,
  keywordstyle=\color{code-kw}\bfseries,
  stringstyle=\color{code-str},
  commentstyle=\color{code-cmt}\itshape,
  numberstyle=\tiny\color{durins-gray},
  numbers=left,
  numbersep=8pt,
  breaklines=true,
  frame=single,
  framerule=0.5pt,
  rulecolor=\color{durins-blue!40},
  xleftmargin=1.5em,
  aboveskip=0.8em,
  belowskip=0.8em,
  tabsize=4,
  showstringspaces=false,
  captionpos=b,
}
\lstdefinestyle{cpp}{
  language=C++,
  backgroundcolor=\color{code-bg},
  basicstyle=\ttfamily\footnotesize,
  keywordstyle=\color{code-kw}\bfseries,
  stringstyle=\color{code-str},
  commentstyle=\color{code-cmt}\itshape,
  numberstyle=\tiny\color{durins-gray},
  numbers=left, numbersep=8pt,
  breaklines=true,
  frame=single, framerule=0.5pt,
  rulecolor=\color{durins-blue!40},
  xleftmargin=1.5em,
  aboveskip=0.8em, belowskip=0.8em,
  tabsize=4, showstringspaces=false,
  captionpos=b,
}
\lstdefinestyle{bash}{
  language=bash,
  backgroundcolor=\color{code-bg},
  basicstyle=\ttfamily\small,
  keywordstyle=\color{code-kw},
  breaklines=true,
  frame=single, framerule=0.5pt,
  rulecolor=\color{durins-blue!40},
  xleftmargin=1.5em,
  numbers=none,
  aboveskip=0.8em, belowskip=0.8em,
  showstringspaces=false,
  captionpos=b,
}

% ── Section formatting ──────────────────────────────────────────
\titleformat{\section}
  {\color{durins-blue}\Large\bfseries}{{\thesection}}{1em}{}[\color{durins-blue}\titlerule]
\titleformat{\subsection}
  {\color{durins-blue!80}\large\bfseries}{{\thesubsection}}{1em}{}
\titleformat{\subsubsection}
  {\color{durins-gray}\normalsize\bfseries}{{\thesubsubsection}}{1em}{}

% ── Header / footer ─────────────────────────────────────────────
\pagestyle{fancy}
\fancyhf{}
\fancyhead[L]{\small\color{durins-gray}Durin's Code — CS4031 Compiler Construction}
\fancyhead[R]{\small\color{durins-gray}\leftmark}
\fancyfoot[C]{\small\color{durins-gray}\thepage}
\renewcommand{\headrulewidth}{0.4pt}
\renewcommand{\footrulewidth}{0.4pt}

% ── Hyperlink style ─────────────────────────────────────────────
\hypersetup{
  colorlinks=true,
  linkcolor=durins-blue,
  urlcolor=durins-blue,
  citecolor=durins-gold,
  pdftitle={Durin's Code — CS4031},
  pdfauthor={Huzaifa Abdul Rehman, Abdul Moiz Hussain, Muhammad Abdullah Khan},
}

% ── Handy boxes ────────────────────────────────────────────────
\newmdenv[
  linecolor=durins-blue!60,
  backgroundcolor=durins-blue!5,
  roundcorner=4pt,
  leftmargin=0, rightmargin=0,
  innerleftmargin=10pt, innerrightmargin=10pt,
  innertopmargin=8pt, innerbottommargin=8pt,
]{infobox}


\begin{document}

% ════════════════════════════════════════════════════════════════
%  TITLE PAGE
% ════════════════════════════════════════════════════════════════
\begin{titlepage}
  \centering
  \vspace*{2cm}
  {\color{durins-blue}\rule{\linewidth}{2pt}}\\[1em]
  {\Huge\bfseries\color{durins-blue} Durin's Code}\\[0.5em]
  {\LARGE\color{durins-gold} Team Reflection Report}\\[1em]
  {\color{durins-blue}\rule{\linewidth}{2pt}}\\[2em]

  {\large CS4031 --- Compiler Construction \\ Spring 2026}\\[3em]

  \begin{tabular}{ll}
    \toprule
    \textbf{Student Name} & \textbf{Student ID} \\
    \midrule
    Huzaifa Abdul Rehman   & 23K-0782 \\
    Abdul Moiz Hussain     & 23K-0553 \\
    Muhammad Abdullah Khan & 23K-0607 \\
    \bottomrule
  \end{tabular}

  \vfill
  {\small\color{durins-gray} Document version 1.0 \quad|\quad \today}
\end{titlepage}

\tableofcontents
\newpage

% ════════════════════════════════════════════════════════════════
\section{Project Summary}
% ════════════════════════════════════════════════════════════════

Durin's Code is a complete compiler for a Tolkien-themed text-adventure
domain-specific language, built from scratch in C++17 as the CS4031 course project
for Spring 2026.
We implemented the full seven-phase compilation pipeline:
a hand-written lexer, a recursive-descent parser producing a typed AST, a two-pass
semantic analyser with a global symbol table, a Three-Address Code (TAC) intermediate
representation generator, a two-pass optimiser (dead-code elimination and redundant-AND
folding), a JSON bytecode serialiser, and a virtual machine with an interactive game loop.

The final system compiles and executes \texttt{.dc} programs entirely from the terminal
with no external runtime dependencies beyond the vendored \texttt{nlohmann/json} v3.11.3
library.
The test suite comprises 49 individual unit tests across six GoogleTest suites, all of
which pass.

% ════════════════════════════════════════════════════════════════
\section{Division of Work}
% ════════════════════════════════════════════════════════════════

\begin{center}
\begin{tabular}{lll}
\toprule
\textbf{Phase / Deliverable} & \textbf{Primary Owner} & \textbf{Student ID} \\
\midrule
Lexer (\texttt{lexer.cpp}) + lexer tests         & Abdul Moiz Hussain     & 23K-0553 \\
Parser (\texttt{parser.cpp}) + parser tests      & Abdul Moiz Hussain     & 23K-0553 \\
Semantic Analyser + semantic tests               & Muhammad Abdullah Khan & 23K-0607 \\
TAC Generator + Optimiser + TAC tests            & Muhammad Abdullah Khan & 23K-0607 \\
Code Generator + codegen tests                   & Huzaifa Abdul Rehman   & 23K-0782 \\
Virtual Machine + VM tests                       & Huzaifa Abdul Rehman   & 23K-0782 \\
CLI (\texttt{main.cpp})                          & Huzaifa Abdul Rehman   & 23K-0782 \\
Language Reference (LaTeX)                       & Muhammad Abdullah Khan & 23K-0607 \\
Architecture Document (LaTeX)                    & Huzaifa Abdul Rehman   & 23K-0782 \\
Example programs (\texttt{examples/})            & All three members      & --- \\
Integration testing + final smoke tests          & All three members      & --- \\
\bottomrule
\end{tabular}
\end{center}

% ════════════════════════════════════════════════════════════════
\section{What Went Well}
% ════════════════════════════════════════════════════════════════

\subsection{Clear Phase Boundaries}
Organising the project as one CMake static library per phase
(\texttt{durins\_lexer}, \texttt{durins\_parser}, and so on) meant each team member
could own a phase independently.
Changes in one phase rarely broke another because the only shared surface was the
typed header interface.
GoogleTest suites gave immediate feedback whenever an interface changed.

\subsection{Memory Safety via \texttt{unique\_ptr}}
Using \texttt{unique\_ptr} throughout the AST eliminated entire classes of memory
bugs from the start.
Ownership semantics were strict --- every child node had exactly one owner --- which
paid dividends when we later added move semantics to \texttt{push\_back} calls.
There were no use-after-free or double-free issues in the parser or downstream passes.

\subsection{Context-Sensitive Lexing for the Dot Operator}
Rather than adding parser lookahead to distinguish \texttt{player.has\_item(\ldots)}
from a generic attribute access, we encoded the distinction directly in the lexer by
tracking \texttt{lastEmitted} and \texttt{lastBeforeDot}.
This kept the parser simpler and produced a richer, more semantically meaningful
token stream.

\subsection{TAC as a Stable Interface}
Once the 18-opcode instruction set was agreed upon, both the TAC generator and the VM
could be developed independently against that contract.
The JSON bytecode format doubled as a clean separation boundary \emph{and} gave us
the \texttt{-o} bytecode-export feature essentially for free.

\subsection{Test-Driven Discovery of Bugs}
Writing test suites before or alongside each phase surfaced several non-obvious bugs
early --- most notably the parser priming bug, the \texttt{unique\_ptr} copy issue,
and the double-line-count in the lexer.
All 49 tests now pass.

% ════════════════════════════════════════════════════════════════
\section{Challenges We Faced}
% ════════════════════════════════════════════════════════════════

\subsection{Parser Priming Bug}
The original parser never called \texttt{advance()} after \texttt{initLexer()}, leaving
\texttt{parser.current} zero-initialised.
Every compilation attempt failed with ``Expect 'room' or 'action' at line 0 col 0''
before reading a single token.
The fix was a single \texttt{advance();} call at the top of \texttt{parseProgram()},
but diagnosing it required careful comparison of the parser's internal model against the
actual token stream state.

\subsection{\texttt{unique\_ptr} Copy vs.\ Move}
Several \texttt{push\_back(ptr)} calls on \texttt{vector\textless unique\_ptr\textless T\textgreater\textgreater}
compiled silently under GCC but were correctly rejected by Clang as illegal copies.
Systematically replacing every \texttt{push\_back(ptr)} with
\texttt{push\_back(std::move(ptr))} resolved the issue and surfaced the correct ownership
semantics.
This bug taught us that relying on one compiler's leniency can hide portability problems.

\subsection{Lexer Line and Column Tracking}
The original \texttt{advance()} function incremented \texttt{lexer.line} on
\texttt{\textbackslash n}, but \texttt{skipWhiteSpace()} also incremented it for the same
newline character, causing double-counting.
Windows \texttt{\textbackslash r\textbackslash n} endings additionally counted as two
separate newlines.
Fixing both required establishing a single rule: \texttt{advance()} is the \emph{only}
function that may modify \texttt{lexer.line} and \texttt{lexer.column}.

\subsection{Two-Pass Semantic Analysis Ordering}
Exits reference rooms that may be declared later in the file.
A naive single-pass approach incorrectly reported ``room not found'' for all forward
references.
Restructuring into two passes --- register everything first, validate second --- resolved
this completely but required careful separation of registration and validation logic into
distinct recursive walks of the AST.

\subsection{Condition Lowering to TAC}
Generating correct TAC for compound conditions (\texttt{\&\&}, \texttt{||}, nested
\texttt{has\_item}) required careful management of temporary boolean registers and
consistent label numbering for \texttt{JUMP\_IF\_FALSE} / \texttt{JUMP} pairs.
Getting label numbering correct for nested \texttt{if} statements --- particularly
when no \texttt{else} branch is present --- required several iterations of the
\texttt{emitCondition()} function.

% ════════════════════════════════════════════════════════════════
\section{What We Would Do Differently}
% ════════════════════════════════════════════════════════════════

\subsection{Integration Tests Between Adjacent Phases}
We wrote per-phase unit test suites but relied on full end-to-end runs to test
the integration between adjacent phases (e.g., whether the semantic analyser
correctly feeds type information to the TAC generator).
Dedicated integration tests at each phase boundary would have caught interface
mismatches earlier and more precisely.

\subsection{Richer Static Type System}
The current type system is implicit: property types are inferred from their literal
values at parse time and not enforced in subsequent phases.
A formal type environment threaded through the pipeline would enable better error
messages for invalid assignments such as \texttt{player.health = "text"}
and would lay the groundwork for type-driven optimisations.

\subsection{Improved REPL Experience}
The \texttt{--interactive} REPL compiles a complete snippet on double-Enter.
A proper incremental REPL would support persistent room and action definitions
across submissions, allowing iterative world-building from the command line.
This would make the tool significantly more useful for rapid prototyping.

\subsection{Error-Tolerant Parser}
The current parser uses panic-mode recovery, which sometimes discards large portions
of the input to find the next synchronisation point.
A more sophisticated error-recovery strategy (e.g., insertion/deletion of a single
token) would produce better error messages for common mistakes like a missing
closing brace.

% ════════════════════════════════════════════════════════════════
\section{Lessons Learned}
% ════════════════════════════════════════════════════════════════

\paragraph{Compiler phases are natural unit boundaries.}
Each phase has a clearly typed input and output.
This makes parallel development practical, makes it easy to substitute
implementations, and produces a codebase that can be read one phase at a time.

\paragraph{Static validation pays dividends at runtime.}
Every semantic check we added to the compiler saved us from writing equivalent
defensive code in the VM.
The VM could be written assuming valid input, making it significantly simpler
and more predictable.

\paragraph{Intermediate representations decouple work.}
The TAC program and JSON bytecode are explicit, inspectable intermediate forms.
The \texttt{--debug} flag that dumps the symbol table and optimised TAC was
invaluable during development.
We learned: build your IR to be human-readable from day one.

\paragraph{Test suites are not optional for compiler development.}
The 49 tests across six suites caught regressions immediately whenever we
modified a shared data structure.
Without them, cross-phase changes would have required manual end-to-end
testing for every modification --- a process that scales poorly.

\paragraph{Ownership semantics should be explicit.}
Using \texttt{unique\_ptr} made ownership immediately visible in the type signature.
By contrast, the raw-pointer approach of the original lexer (which stores a pointer
into the source buffer) required careful documentation to be safe.
The lesson: prefer types that express ownership invariants over comments that describe
them.

% ════════════════════════════════════════════════════════════════
\section{Conclusion}
% ════════════════════════════════════════════════════════════════

Building Durin's Code from a proposal to a working compiler with 49 passing tests
gave our team practical experience with every stage of the compilation pipeline.
The project demonstrated that a relatively small DSL --- fewer than 15 keywords
and a handful of statement types --- nevertheless requires a carefully engineered
multi-phase compiler to be robust and maintainable.

The challenges we encountered (parser priming, ownership semantics, two-pass analysis,
condition lowering) are the same classes of problems faced in production compiler
development.
Working through them ourselves, rather than studying them in the abstract, produced
a significantly deeper understanding of the material.

\end{document}
